% ============================================================
%  CHAPTER VI
%  Effects of Reverse Perspective
%  Source pp. 118–139
% ============================================================

\chapter{Effects of Reverse Perspective}
\markboth{Chapter VI. Effects of Reverse Perspective}{Chapter VI}

\srcpage{118}

If medieval art contained only those features of perspective construction that were
discussed in the preceding chapter, the analysis of the geometry of works of painting
would be a comparatively simple task.  In reality the question is considerably more
complex.

When one attempts to depict near regions of space, the conclusion set out in
Chapter~III makes itself felt with full force: the transmission of the geometry of
visual perception of near space without apparent departures from that geometry
(``distortions'') is in principle impossible, regardless of which perspective system
is used.

Naturally, the medieval master had no inkling of such a theorem, but he felt it in
practice.  Artists knew this and passed their knowledge from master to pupil.  Such
knowledge existed not in an abstract theoretical form but was transmitted by
instruction in traditional modes of depiction.  The pupil was shown what and how it
was most reasonable to distort in each case, and what to render correctly.  Of
course, depending on the culture and its traditions such instructions might differ.
The range of techniques was very wide, as is apparent in the medieval visual arts of
different regions.

The principled impossibility of transferring the visible geometry of near regions of
space onto a plane, and the circumstances connected with this, nevertheless left the
medieval master great freedom.  Moreover, he felt himself master of the situation,
having the right to choose.  This ``right to distort,'' arising from the then-unknown
but objectively existing geometrical laws, gave rise to a feeling of subjective
freedom.  If the artist himself had to freely introduce distortions, he introduced
them where and to the degree that suited his artistic intentions~--- and, let us add,
sometimes quite independently of whether or not geometry required them.  For if
distortions are unavoidable, let them at least be artistically justified and even
artistically necessary.

\srcpage{119}

On this path into the artistic practice of painters there enter not only geometrically
conditioned distortions of visible images (such as those examined in the preceding
chapter in connection with the problem of the form-constancy mechanism), but also
distortions utterly unconnected with any attempt to accurately transmit visual images
on a plane.  Consequently, an important result of the action of the geometrical laws
mentioned above was the possibility of departing from the ``original,'' i.e.\ the
possibility of not following the visual image and of transforming the depiction with
the aim of enhancing its expressiveness.

The transformations referred to here altered the geometrical properties of images and
could therefore influence the formal characteristics of perspective.  Since these
transformations were not the result of an attempt to convey the geometry of visual
perception as accurately as possible, but had other causes, it is appropriate to
speak here of perspective \emph{effects}.  One or another alteration of perspective
constructions conditioned by such transformations is yet another consequence of
artistically justified impulses.  Naturally, primary attention will be devoted to
those impulses that led to effects of reverse perspective.  Below the following
questions will be examined (preserving the numbering begun in Chapter~V): 5~---
allowance for the mobility of the viewpoint; 6~--- the striving to increase
informativeness; 7~--- the requirement of ``non-occlusion''; 8~--- hierarchical
considerations; 9~--- compositional requirements; and 10~--- allowance for the
functional peculiarities of painting.

\medskip
\noindent\textbf{5.}\quad
\emph{Allowance for the mobility of the viewpoint.}  The cases to be studied in this
section can only partially be classed among those leading to perspective
transformations.  They might equally be assigned to tasks connected with the wish to
convey visual impressions more accurately.  These are, in a certain sense,
intermediate phenomena, and they are assigned to the perspective-effect-producing
group only because this is predominantly their character.

Let us begin with the effects of direct perspective.  When in an interior, a person
immediately becomes aware that he is unable to see it with equal sharpness in all
directions simultaneously.  To gain a sufficiently detailed impression of it, it
proves necessary to shift the gaze successively from one part to another.  If one
additionally supposes that the room in which the person stands is not too large, then
wherever he directs his gaze, within the narrow field of sharp vision he will see the
axonometric structure of that section of perceptual space.

\srcpage{120}

As already noted above, axonometry is the perceptual perspective of near regions of
space only in a first approximation; but for simplicity in what follows we shall
confine ourselves in this chapter to the first approximation~--- the more so since
artists too often confine themselves to it.  The conclusions reached below will then
be highly clear, and refinements of the kind examined in the preceding chapter can
always be added without difficulty.  Moreover, this approach is well consistent with
the notion introduced above of axonometry as the perspective basis of medieval
painting.

If the artist, successively shifting his gaze, attempts to depict everything exactly
as he sees it, his picture will ``be composed'' of a series of local axonometries.
The word ``composed'' is placed in quotation marks in part because such a mechanical
composition cannot in fact be carried out: local axonometries will come into
screaming contradiction with one another (see Ill.\,23).

Before turning to a discussion of techniques capable of ``smoothing out'' the obvious
absurdities when reconciling local axonometries, one should draw attention to one
extremely important circumstance.  Everything written above holds equally for the case
when the person in the room stands still.  Wishing to form a more complete idea of
the interior, he may do so by moving about the room.  It is easy to see that in this
case the left wall of the room and the space adjoining it will be especially well seen
when the person stands on the right, and vice versa.  If the artist is not bound by
the requirement of depicting the interior from a single viewpoint, he will almost
certainly prefer a depiction characterised by a mobile viewpoint.  The difficulties
of reconciling local axonometries will increase, although the nature of those
difficulties will not change.

As already noted in Chapter~III (Ill.\,18), axonometry does not allow the interior
to be depicted from within; this is possible only by partial depiction, for example
of one corner (scheme~B).  This circumstance pushed the antique artist toward using
several viewpoints for the simultaneous depiction of all four corners of an interior,
each of which can be depicted according to scheme~B of Ill.\,18 if the viewpoint for
each corner is correctly chosen.  Using several viewpoints was considered admissible
even in the Renaissance and was entirely natural in antiquity, which did not know the
strict theory of linear perspective.

The antique artist's desire to depict the right part of an interior from the left and
the left part from the right was entirely natural.  The problems arising with such a
mode of depiction can be understood by turning to the scheme in Ill.\,45.  The left
part of the figure gives four parallelepipeds that visually represent the types of
axonometries most natural for depicting each of the four corners.  Each of these
axonometries corresponds to its own viewpoint and direction of gaze, and the edges
forming the axonometric structure of the corresponding corner are shown in thick
lines.  The right part of the figure gives the depiction of the interior by means of
these four independent axonometries, with thick lines showing the edges (junctions of
floor, walls, ceiling) parallel to those shown in the source axonometric
parallelepipeds.  As is visible from the schemes in the left part of the figure, the
upper right corner is given when seen from the left and below, the lower right corner
when seen from the left and above, and so on.  Subsequently each corner of the
interior is given in its own axonometry (from its own viewpoint).  From the right
part of the figure it is visible that if the rectilinear structure of the depicted
surface is parallel to the axes of neighbouring axonometries having the same
direction, then the reconciliation of such structures of two adjacent corners presents
no difficulty.  Indeed, if one needs to depict the vertical structure of the walls
(for example, columns, etc.), they will be depicted as vertical lines in both the
lower and upper axonometries and will join without difficulty.

\srcpage{121}

If the structure of a surface is parallel to the axes of neighbouring axonometries
that do not have the same direction in the figure, a principled impossibility of
contradiction-free reconciliation of adjacent images arises.  In Ill.\,45 this is
shown by the thick triangle that can appear when depicting coffeered ceilings
(sometimes this will be not a triangle but a trapezoid).\footnote{In examining the
  question of the mobility of the viewpoint in antique and medieval art, the analysis
  of the oft-described method of depicting the architectural background placed in the
  second plane~--- in which each building and even its individual parts were often
  depicted from different viewpoints~--- is deliberately omitted here.  Below only
  objects in the foreground will be discussed.}

The incompatibility of reconciling the left and right axonometries amounts to this:
in both the left and right corners, groups of objectively parallel lines are parallel
in the figure as well, while in the central part of the figure two such lines turn
out to be disposed at some angle to each other.  To avoid having to depict figures of
this kind, which undoubtedly spoil the picture by introducing a certain dissonance,
one may proceed in different ways.  The simplest is not to depict parallel structures
of surfaces at all.  From the example of depicting the floor in Ill.\,45 it is
visible that since the left and right walls of the interior are separated by a
sufficiently large distance, the non-parallelism of the images of the floor-wall
junctions does not catch the eye.  This is consistent with the artists' repeatedly
mentioned striving in antiquity and especially the Middle Ages to depict as parallel
those lines that are not only parallel to each other in real space but are also
separated from one another by a small distance.

Another way to avoid having to depict the unnatural triangle or trapezoid in parallel
structures is to conceal this location by a break, a drapery, or some other means.
On a Pompeian fresco of Style~IV, a coffeered ceiling corresponding to the described
scheme is depicted (Ill.\,46).  Careful analysis shows that the lower ceiling is
given in two strict axonometries (this is easy to do, since the mutual inconsistency
of the left and right parts is removed by a break in the ceiling in the middle), while
the linear structure of the upper ceiling is shown not by parallel lines but by
slightly converging straight lines.\footnote{The striving to conceal the impossibility
  of reconciling the left and right axonometries is observed not only when depicting
  ceilings.  Thus, for example, the roof of the palace of Theodoric in the mosaic of
  the basilica of Sant'Apollinare Nuovo in Ravenna (6th century) is depicted by means
  of three systems of parallel lines~--- inclined in different directions for the left
  and right parts of the roof~--- whose mutual inconsistency is removed by the
  depiction of the central roof section shown by a system of vertical lines.  See, for
  example: V.\,N.~Lazarev.  \emph{Vizantiyskaya zhivopis'} [Byzantine Painting].
  Moscow, 1971, p.~53.}
This is interesting in two respects.  First, it is obvious that the weak
transformation of the axonometry was made so that the central trapezoid would become
less noticeable while the linear structure of the ceiling would continue to appear in
each corner as composed of parallel lines.  The artist clearly wants the viewer to
continue to consider that each corner is given in its own axonometry (as in the
depiction of the lower ceiling).  Second, it is obvious that the artist is unaware of
linear perspective; otherwise, by increasing the convergence of the lines forming the
structure of the upper ceiling, he could have rid himself of the unnatural trapezoid
in the centre of the image.

\figblock{45}{Scheme for depicting an interior using local axonometries.}

\figblock{46}{Pompeian fresco of Style~IV from the house of Epidius Sabinus.
  Detail (tracing).}

\srcpage{122}

Even if linear perspective is unknown to the fresco's author, the combined action of
several local axonometries (four in Ill.\,45) creates an integral scheme of interior
depiction in direct perspective: the width and height of the space decrease in the
picture as one moves away from the viewer.  Thus the use of several viewpoints, while
maintaining local axonometry as the basis of individual parts of the picture, leads
in the case examined to the overall appearance of direct perspective.

Medieval art too encountered the problems of reconciling local axonometries when
depicting interiors.  As an example let us cite the predella by the Italian painter
of the second half of the 14th century (Barnaba da Modena, ``Scenes from the Life of
Christ'') on which is depicted a composition on the theme of the ``Presentation in
the Temple'' (Ill.\,47).  The left and right parts of the ceiling are rendered in the
system of local axonometries, and the artist concealed their mutual inconsistency by
showing in the middle part

\srcpage{123}

of the ceiling a hexagonal recess~--- a kind of baldachin placed directly over the
altar.  This compelled giving the altar too an unusual form based on a hexagon.  Such
a composition could solve the problem only partially, and in order visually to justify
the hexagonal recess in the ceiling, the capitals of the very slender columns shown
in the foreground carry certain structural elements of unclear purpose, which
nonetheless turn out to be parallel to the corresponding ceiling structures (this is
especially true of the left column; at the right one the artist was constrained by
the boundary of the adjacent image).

It should be emphasised that the whole problem was solved only visually.  If one
attempts to pass to the objective space shown in this predella, then since the left
and right parts of the ceiling are composed of parallel boards having the same
direction in objective space, the corresponding sides of the hexagon in the central
part of the ceiling must also be parallel to each other.  But then the depicted
figure cannot be a regular hexagon, in which, as is well known, these sides are
inclined to each other at 60°.  Consequently the depicted spatial formations are
formally geometrically absurd, yet they are composed so that they are not only not
perceived as absurd but, on the contrary, lend the whole image geometric conviction.

\figblock{47}{Barnaba da Modena.  ``The Presentation in the Temple.''  Second half of
  the 14th century (tracing).  Museum of the Trinity-Sergius Monastery (MZhVI).}

Suppose now that it is necessary to show an object, some group, etc., from
``outside'' rather than ``inside,'' under the condition that the left and right parts
of the depicted scene are equally important.  Then it is natural to use two viewpoints
and to depict the left part from the left and the right part from the right.  If each
part is depicted in its own axonometry, the result of their ``gluing'' will be an
image in strong reverse perspective (Ill.\,48).  An example of such strong reverse
perspective as a result of using two viewpoints is the icon ``New Testament Trinity''
from the collection of the Zagorsk Museum (Ill.\,49).  The reverse perspective of the
throne is of the origin described above; this is attested in particular by the fact
that each of the two sides of the throne, considered separately, as well as the
footstool, is given in axonometry.  Here the desire to show the Father and the Son in
exactly the same way, without mutual occlusions or any otherwise (perhaps
unintentionally) expressed ``inequality of substance'' of the depicted persons,
manifests itself obviously.  The ``gluing'' of the throne may also be interpreted as
a symbol of the indivisibility of the Trinity, i.e.\ it may have not only a
compositional significance.  It is interesting to note that the front (and back)
edges of the ``glued'' seats are not parallel, with the result that the ``gluing'' is
carried out by curved lines.  The latter indicates that here there is no simple
depiction of a single throne in reverse perspective, but a conditional combined
depiction of two seats turned in space in accordance with the poses of the Father and
the Son.

\figblock{48}{Depiction of a parallelepiped from two viewpoints by ``gluing'' two
  axonometric images.}

\srcpage{124}

\figblock{49}{``New Testament Trinity.''  Icon of the first half of the 16th century.
  Zagorsk Museum of History and Art (ZIKhMZ).}

\srcpage{125}

Sometimes the fact of ``gluing'' two axonometries is completely obvious.  In this
respect the depiction of the Evangelist John dictating to Prokhor (Ill.\,50), placed
in a carved frame of royal doors from the church of the village of Sofrino (17th
century), is very instructive.  Prokhor is painted on the right and John on the left,
each on his own seat; their feet are brought so close together that it becomes
impossible to give each his own footstool.  The artist combines these two footstools,
``gluing'' them, and as a result this single footstool is depicted in the strongest
possible reverse perspective, clearly contradicting the whole geometrical structure of
the icon.  This is, of course, not a perspective technique but a perspective effect
arising from compositional impulses (conditioned by the mobility of the viewpoint).

\figblock{50}{``St.~John the Theologian and Prokhor on the Island of Patmos.''
  Compartment of royal doors from the church of the village of Sofrino.  17th century.
  Tracing.  State Museum-Reserve of the Moscow Kremlin (GMZKr).}

What has been said in this section allows one to assert that the mobility of the
viewpoint leads to various transformations of the image.  If the basis of the picture
is local axonometry, then the use of several viewpoints can lead both to direct and
to reverse perspective in the overall image.  The examples given and elementary
considerations lead in particular to the conclusion that the cumulative effect of
direct perspective arises when the artist, as it were, ``crosses'' the directions of
the projecting lines~--- i.e.\ observes the right parts from the left, the left from
the right, the upper from below, and so on.  In the opposite case~--- when no such
crossing occurs and the right parts are observed from the right, the left from the
left, and so on~--- the cumulative effect of reverse perspective arises.  The first of
these cases is especially characteristic of depictions from ``inside'' (e.g.\
interiors), the second for depictions of not-too-large objects, groups, etc., from
``outside.''

\srcpage{126}

The predominant appearance in Old Russian and Byzantine painting of effects of the
second type is connected with the fact that the medieval artist did not strive toward
the depiction of interiors,\footnote{When this proved necessary, he preferred, as is
  well known, the more informative depiction of the exterior of the building against
  the background of which the event happening inside it was shown.}
devoting primary attention to depicting saints and the objects associated with them
(seats, footstools, etc.), whereas the antique artist not infrequently depicted
``empty'' interiors, for example as decorative wall paintings.  Thus the manifestation
in one case of a tendency toward reverse perspective and in another toward direct
perspective (while preserving local axonometry) is connected with what was depicted,
since in a certain sense this also determined how it was depicted.

The multiplicity of viewpoints is known in later art as well.  One need only recall
the examples of pictures by Raphael (``School of Athens''), Paolo Veronese (``The
Wedding at Cana,'' ``The Feast at the House of Levi''), and other artists.  However,
in later art the motivating cause of such a multiplication of viewpoints is, strange
as it may sound, the desire for a more accurate conveyance of visual perception from
a single viewpoint.  If the viewpoint is shifted vertically~--- the floor painted from
a high viewpoint and the ceiling from a low one~--- the vanishing point of the floor
will rise and that of the ceiling will fall.  But then within the depicted space the
convergence of objectively parallel lines of both floor and ceiling will be (compared
with linear perspective) weakened, and their depiction will be brought closer to the
visual perception (perceptual perspective) corresponding to a single viewpoint
situated somewhere between these two.  It should be recalled that compensations of
this kind in the system of perceptual perspective are connected with the enhancement
of distortions of other elements of the picture~--- in this case, the walls.  For
them the convergence will not be weakened but enhanced.  Therefore, when it is not the
floor and ceiling but the walls that are important, the floor must be painted from a
low viewpoint and the ceiling from a high one, which will lead to improved depiction
of the walls at the cost of worsened depiction of the floor and ceiling.

\medskip
\noindent\textbf{6.}\quad
\emph{The striving to increase the informativeness of the picture.}  As is well
known, one of the aims of medieval painting was to communicate the maximum information
about the depicted object to the viewer, even at the cost of violating perspectival
correctness.  An increase in informativeness can in particular be achieved by
depicting individual parts of an object turned so that they can be better ``made
out.''

\srcpage{127}

In the depiction of the writing monk Ugone (Ill.\,51)~--- which is not a saint's
image at all but closest to a portrait~--- the tilt of the writing desk is so steep
that here the desire to show the viewer the very process of writing in all its detail
is clearly at work.  The tilt arising from such informational impulses produced an
effect of reverse perspective (the far end of the desk is higher than the near end)~---
the more remarkable in that the desk itself is given in this Italian 14th-century
fresco in direct perspective.  This rotation of the desk board can, if one wishes,
also be interpreted as the result of a mobile viewpoint~--- to show the details of
writing one must raise the viewpoint.

\figblock{51}{Tommaso da Modena.  ``Brother Ugone de Provence.''  1352.  Treviso,
  Santo Niccolò, chapter room.  Detail of fresco (tracing).}

This last remark leads to a quite natural question: is it in general possible to
distinguish mobility of viewpoint from artificial rotations of depicted surfaces?  In
pure geometry both effects are indistinguishable, and so here the peculiarities of
the picture as a whole must be taken into account.  As a rule, a mobile viewpoint
leads to a transformation of the whole ensemble of observed objects; for example, in
the icon ``New Testament Trinity'' from the Zagorsk Museum, not only the seat but
also the figures are painted from two viewpoints: Christ is depicted entirely from
one and the Father from the other.  With an informational rotation of a certain
surface, it is quite admissible for everything else to undergo no similar deformation.

One should always remember that clear-cut classifications in art~--- which admits
gradual transitions and the simultaneous action of different causes~--- can only be
conventional.

Old Russian and Byzantine painting is full of examples of depictions of surfaces that
are clearly given artificial rotations of informational significance.  Not
infrequently, in order to allow the viewer who has come close to an icon not only to
see the image of an open Gospel book but to read what is written on its pages, the
artist ``rotates'' the surface~--- for example of the small table on which the Gospel
lies~--- in the same way as the desk board is rotated in Ill.\,51.

In general, informational rotations of horizontal surfaces~--- to show in detail the
objects arranged on them~--- are very frequent, especially when tables with food,
vessels, and so on are depicted on an icon or fresco.

\srcpage{128}

As an example one may point to the table in the compartment ``Nicholas returns Vasily
to his parents'' of the icon ``Nicholas of Zaraisk with his Life'' (Ill.\,53).

In all these and many analogous cases of depictions of horizontal surfaces about which
enhanced information is to be given, the effect of reverse perspective always arises,
since all such surfaces are depicted when viewed from the front and above.  If the
need were to show them from the front and below, a rotation of the surface for a
detailed depiction of its underside would lead to a reduction of the back legs~--- i.e.\
to the effect of direct perspective.  Thus informational rotations of horizontal
surfaces lead to effects of reverse perspective only insofar as one must practically
always increase information about the upper sides of such surfaces.

\figblock{52}{``Nicholas leads Vasily from Saracen captivity.''  Compartment of the
  icon ``Nicholas of Zaraisk with his Life.''  16th century.  State Gallery
  of Fine Arts (GGKhM).}

\figblock{53}{``Nicholas returns Vasily to his parents.''  Compartment of the icon
  ``Nicholas of Zaraisk with his Life.''  16th century.  State Gallery
  of Fine Arts (GGKhM).}

\srcpage{129}

The informational rotations discussed here act in the direction of the form-constancy
mechanism examined in detail in this chapter.  Perhaps this is why artists readily
used this technique: it merged both with the effects produced by the form-constancy
mechanism and with the modes of depicting near objects in natural reverse perspective.
The rotation of a depicted horizontal surface for reasons of increasing information
is, as a rule, of noticeably greater magnitude than the rotation produced by the
form-constancy mechanism alone.

Informational rotations of vertical planes can also lead to corresponding
transformations of perspective character.  As examples let us refer to the schemes
(Ill.\,54) representing the perspective constructions that arose in the miniature of
the 13th-century Gospel ``The Healing of the Paralytic'' and in the miniature from
the Menology of Basil~II, ``Theodore Stratelates.''  In the first case the artist
wished to show simultaneously the inner lateral slopes of the roofs of two parallel
structures directed ``away from the viewer,'' as a result of which an effect of direct
perspective arose; in the second case the two outer lateral walls of a building, also
directed away from the viewer, are shown simultaneously, as a result of which an
effect of reverse perspective arose~--- the more remarkable in that the side walls of
the building are shown decreasing with depth.\footnote{For the cited miniatures see,
  for example: V.\,N.~Lazarev.  \emph{Istoriya vizantiyskoy zhivopisi} [A History of
  Byzantine Painting], vol.\,I.  Moscow, 1948, table~X\,b; vol.\,II, table~71.}

\srcpage{130}

\figblock{54}{Appearance of the effects of direct and reverse perspective as a result
  of informational rotations of vertical planes.}

If in the depiction of horizontal planes the desire to increase information about them
practically always leads to reverse perspective, this cannot be asserted for the
depiction of vertical surfaces.  However, rotations of vertical surfaces occur much
less often than rotations of horizontal ones, and even vertical planes themselves upon
rotation in half the cases (and actually more often) also give an effect of reverse
perspective.  Consequently the general result is that all kinds of artificial rotations
of the planes of depicted objects with the aim of increasing information about the
object in the vast majority of cases lead to images formally possessing the signs of
reverse perspective.

\medskip
\noindent\textbf{7.}\quad
\emph{The ``non-occlusion'' requirement.}  In medieval art one frequently finds the
artist's striving not to allow one object or figure to be occluded by others.

\srcpage{131}

The reasons for this are usually unconnected with the modes of spatial construction.
This approach can be explained by the desire not merely to depict what is visible but
to narrate the event, ``arranging'' everything needed for the narrative so that one
thing does not ``interfere'' with another.  The reasons for such a compositional
approach cannot be examined in more detail here, but since such a tendency exists, it
can affect the geometry of the image and accordingly lead to perspective effects.

If one turns again to the fragment of the icon ``The Presentation in the Temple''
(Ill.\,31), one sees how the artist's striving not to let the image of Anna's garment
and Mary's footstool intersect leads to a clear reverse perspective of the footstool.
At the same time the adjacent steps on which Simeon stands the artist is not inclined
to paint in an analogous way.  Indeed, the reverse perspective of the footstool has a
magnitude of 26°, while the lower step is rendered in slight direct perspective (6°)
and the upper step in slight reverse perspective (13°).  The impression is created
that the artist painted everything ``on average'' axonometrically, making an exception
only for Mary's footstool.

In the icon ``The Dormition'' from the Kirillo-Belozersky Monastery, an angel cutting
off the hands of the impious Avfoniy is depicted before the bier of the Mother of
God.  If this group were given on the same scale as the other figures (apostles,
bishops, women), the angel and Avfoniy would occlude the image of the Mother of God.
Precisely for this reason they are shown at a reduced scale, scarcely above the knees
of the apostles (Ill.\,69).

Even more compelling compositions with the same motivating impulse can be cited from
the medieval world beyond Old Russian and Byzantine art.  In the two examples given
above, the appearance of the reverse-perspective effect was connected with local
alterations in the form of an individual object or the scale of an individual figure.
But if one turns, for example, to medieval Western art, one can cite an example of a
composition where, from the requirement of ``non-occlusion,'' a progressive increase
in the scale of human figures with depth is used.  The Carolingian-era miniature
``The Presentation of the Codex to Charles the Bald'' (Ill.\,55) possesses precisely
this feature.  The reduced depiction of the figures in the foreground has the aim of
``freeing up space'' in front of Charles the Bald's throne.

\figblock{55}{``The Presentation of the Codex to Charles the Bald.''  Miniature from
  the dedicatory leaf of the Vivian Bible.  Mid-9th century.  Paris,
  Bibliothèque nationale.}

Examples of this kind are more numerous than might be supposed, and are very widely
distributed.

\medskip
\noindent\textbf{8.}\quad
\emph{Hierarchical considerations} leading to an enlarged depiction of the principal
figures can also produce perspective effects.  If a figure depicted closer to the
viewer is shown smaller than a figure farther away, owing to the higher position of
the distant figure in some hierarchy, an effect of reverse perspective will arise.
True, this will also be rather a formal reverse perspective, since it usually concerns
only human figures.  As an example of such ``hierarchical'' perspective one may point
to the royal doors of 1425--1427 of the Trinity Cathedral of the Trinity-Sergius
Monastery, where the Evangelist John is shown seated farther away than his disciple
Prokhor, yet the figure of the Evangelist is given at an enlarged scale relative to
Prokhor (Ill.\,56).

\srcpage{132}

\figblock{56}{``St.~John the Theologian and Prokhor on the Island of Patmos.''
  Compartment of royal doors (tracing).  1425--1427.  Zagorsk Museum of History
  and Art (ZIKhMZ).}

The number of similar examples is very easily multiplied; this is not done here since
the type of perspective effects described is completely obvious.

\medskip
\noindent\textbf{9.}\quad
\emph{Compositional requirements.}  The causes of perspective transformations examined
in the present chapter show that on the one hand they were conditioned by the
indepictability of the geometry of perceptual space on a plane and from this, on the
other hand, by the medieval master's ``sense of freedom.''  The artist altered spatial
constructions also in cases where the composition demanded it.  Examples have already
been given above that may be assigned to cases of this kind~--- the ``gluing'' of the
footstools of John and Prokhor into one, the reduction of the figures of the angel
and Avfoniy on the icon ``The Dormition,'' and the transformation of the form of the
footstool on the icon ``The Presentation in the Temple.''

The influence of compositional requirements on spatial constructions was much more
varied than follows from the cases mentioned.  To illustrate this let us turn to the
following examples.  In two compartments of the already-mentioned icon ``Nicholas of
Zaraisk with his Life'' (``Nicholas leads Vasily from Saracen captivity'' and
``Nicholas returns Vasily to his parents'') covered tables are depicted.  In both
cases the front edges of the tables are at the same level from the floor, but in one
case the visible surface of the covered table amounts, as a fraction of the front-leg
height, to $0.48$, and in the other to $0.97$.  Such a large difference is connected
neither with the laws of visual perception nor with the desire to give quantitatively
different information about the objects arranged on the tables.  This difference has
a compositional basis~--- in the first case seated figures and in the second standing
figures are shown behind the table.  The artist depicts one and the same object in
different ways, i.e.\ uses his right of transformation not at all solely in order to
seek the best way of depicting the table as such, but in order best to compose the
picture as a whole (Ills.\,52, 53).

\srcpage{133}

In Ill.\,57 is shown the symbol of the Evangelist Matthew from the ``Khitrovo
Gospel.''  The silhouette of the Gospel book differs from the usual ones, in
particular those shown in Ill.\,43.  The lightness of the angel's movement is
emphasised by the sharply acute lower right angle of the Gospel book's silhouette,
even though from perspectival considerations it should have been made obtuse, adding
one more angle to the book's silhouette.  It is completely obvious that the
transformation of the visible form of the book has no perspective nature; nonetheless,
having arisen, it gave the formal effect

\srcpage{134}

of reverse perspective~--- the back board of the binding became wider than the
front.\footnote{One should not think that the features of depicting Gospel books
  described here and above are a specific property of medieval (in particular
  religious) painting.  The same characteristic features are also known in ancient
  Roman painting (for example, the portrait of a Roman woman from the house of
  Libanius at Pompeii: see A.~Chubova.  \emph{Drevnerimskaya zhivopis'} [Ancient
  Roman Painting].  Leningrad\,--\,Moscow, 1966, ill.~42).}

\figblock{57}{Symbol of the Evangelist Matthew.  Miniature from the ``Khitrovo
  Gospel.''  End of the 14th~-- beginning of the 15th century.  State Lenin Library
  (GBL).}

In the miniature with the Evangelist Mark (Ill.\,58) all objects are subject to
``normal'' reverse perspective, except for the seating surface, which is given in a
very strong widening.  This is easily explained by the artist's desire for a clear
silhouette~--- a single smooth line bounding the leg of the seat, its surface, the
back, and the evangelist's head.  Consequently in this case too the strong reverse
perspective of the seating surface arose as a result of compositional rather than
perspective stimuli.

\figblock{58}{Evangelist Mark.  Miniature from a 14th-century Gospel.  State
  Historical Museum (GIM).}

\srcpage{135}

The examples given show that the medieval artist considered himself entitled to
distort the visible forms of objects when this seemed expedient for compositional
reasons.  One can adduce especially many observations of this kind from the
depictions of the four evangelists on royal doors~--- often one and the same author
gives the contours of the footstools of each of the four evangelists a completely
different character, even though there could be no reasons for this other than
compositional ones.

The examples given are sufficient to form an impression of the appearance of reverse
perspective from compositional considerations.  The question is perfectly legitimate,
however, whether compositional considerations might not shift axonometry in the
direction of direct perspective.  Reasoning abstractly, one should expect equal
probability of transformations of the axonometrically based image in both directions.
Confirmation of the possibility of direct perspective appearing in the depiction of
individual objects is provided by the Novgorod icon ``Pokrov,''\footnote{See:
  V.\,I.~Antonova.  \emph{Drevnerusskoye iskusstvo v sobranii Pavla Korina} [Old
  Russian Art in the Collection of Pavel Korin].  Moscow, 1966, ill.~7.}
where, proceeding from the mentioned compositional requirement of ``non-intersection,''
Roman's footstool is made converging with depth, so as not to merge on the right with
the image of the rounded ambon.

Thus, shifts are possible both toward direct and toward reverse perspective.  However,
if one attempts to make a numerical assessment of the number of distortions from the
true form in both directions, an enormous preponderance of transformations bearing the
formal signs of reverse perspective immediately emerges.  Evidently the medieval artist
preferred shifts toward reverse perspective because axonometry of the close foreground
was being transformed in this direction by the whole totality of the causes examined
above.  One may therefore suppose that the artists' choice was guided in the cases
described here by a sense of stylistic unity.

\medskip
\noindent\textbf{10.}\quad
\emph{Allowance for the functional peculiarities of painting connected with its
purpose} affects not only the depiction of individual objects but also the geometric
peculiarities of the construction of the whole space.  The religious painting of the
Middle Ages is distinguished by an enhanced suggestive power, achieved by a variety of
means, among which the geometric peculiarities of spatial rendering (including those
giving effects of reverse perspective) play no small role.  This question will be
examined at the end of Chapter~IX.

\medskip
\noindent
$* \quad * \quad *$
\medskip

\srcpage{136}

The causes examined in the last two chapters of the appearance in medieval painting of
images constructed in the reverse-perspective manner are very varied.  They may be
gathered into three groups.  First, it is the consequence of the undistorted
transmission of visual perception of near regions of space.  Second, it is the
consequence of the fact that volumetric bodies transformed by the constancy mechanisms
cannot be depicted without distortion, and the attempt to transmit the ``main'' thing
undistorted can lead to quite strong reverse perspective.  And finally, third, it is
the consequence of incidental perspective effects resulting from transformations of
images in connection with other goals, usually more immediately dependent on the
artistic nature of the image.  To this last group one should assign the use of a
system of local axonometries in constructing the picture as a whole, the striving to
increase informativeness, and certain compositional considerations.  One should recall
that in the cases enumerated here effects of reverse perspective arise only under
certain conditions~--- in depicting convex bodies from outside rather than interiors,
and in the case of depicting objects when viewed from above and in front.  It is
precisely this that is characteristic of Byzantine and Old Russian art.  To this same
group one should assign those cases where effects of reverse perspective proved to be
the consequence of the compositional requirement of ``non-occlusion,'' the use of
hierarchical differences of scale, and so on.

\srcpage{137}

The multiplicity of causes for one and the same result~--- which today we call
``reverse perspective''~--- led on the one hand to this type of perspective
construction becoming the most conspicuous feature of icon painting, and on the other
to it being practically impossible today to distinguish the impulses.  For example,
looking at the depiction of the tables in Ills.\,52 and~53, it is very difficult to
say to what degree the tilt of their horizontal surfaces was caused by the influence
of the form-constancy mechanism, to what degree by the striving to increase
informativeness, and to what degree by compositional considerations.

Since all the heterogeneous causes examined above acted in the same direction, a
characteristic style arose that led to depicting in reverse perspective even those
objects for which axonometry would more often be natural~--- and that exerted its
influence on the composition of the picture as a whole.  Having arisen, this style
extended to distant planes as well, for which according to the laws of the psychology
of visual perception the rules of direct perspective~--- or, for small depths,
axonometry~--- should have been adhered to.  One cannot assert that the Old Russian
artist did not feel that the near and distant planes must be depicted differently~---
as a rule, the distant planes deviate from axonometry much less markedly than the near
one.  And yet the architectural \emph{staffage} is characterised not only by a
sharply pronounced variability of viewpoint, but also by the appearance in the
depiction of individual buildings of elements of reverse perspective.

\srcpage{138}

In conclusion, some preliminary thoughts on the problem of constructing the space of
the picture as a whole should be given.  This question cannot be reduced to geometry;
the space of a picture is constructed by the whole totality of means available to the
artist.  It nonetheless seems reasonable to extract the geometrical aspect of this
problem, which may contribute to its solution in a more general context.

The basic feature of all the types of spatial construction examined above is
\emph{locality}.  All these considerations concerned the depiction of individual
objects~--- and not too large ones, fitting entirely within the field of sharp
vision~--- rather than space as a whole.  In those cases where it was necessary to
depict a large space (for example, an antique interior), it was as it were divided
into several small sections, each depicted as a local formation, and the ``gluing'' of
these local images proved a difficult (though by no means hopeless) task.  The
psychophysiological basis of the system of local images is the smallness of the human
field of sharp vision.

This feature of the geometry of medieval (and earlier) pictures was often emphasised
in art-historical works when it was noted

\srcpage{139}

that space in art of this type is constructed by the depiction of individual objects.
The geometrical essence of this feature of painting follows most clearly from the
answer to the question of what rules govern the depiction of the ``empty'' spaces
between objects.  In essence, this problem does not preoccupy the antique or medieval
artist.  He finds nothing strange in the fact that if two adjacent rectangular
footstools are shown in reverse perspective, the gap between them will itself turn out
to be in direct perspective.

The system of linear perspective gave a flawless geometrical solution to the problem
of constructing the space of the picture, providing a strict method for depicting any
individual point of pictorial space on the picture plane, as a result of which both
objects and the spaces between them are depicted by uniform rules.  In this connection
two questions naturally arise: how should the space of the picture as a whole be
constructed if the basis of its geometry is the principle of local depictions of
objects; and can a system of perceptual perspective exist for near space in which
objects and the spaces between them are constructed by strict uniform rules?

The first of these two questions is discussed at the end of Chapter~IX; the second is
the subject of Chapter~X of the present book.
